% =====================================================================
%  CHAPITRE 14 — PARLER : SITUATIONS ET DIALOGUES
% =====================================================================
\chapter{40 situations : dialogues, blocs et jeux de rôle}

Ce chapitre est le c\oe ur pratique du livre. Chaque situation donne :
\textbf{(1)} un dialogue complet, traduit et prononcé ; \textbf{(2)} les blocs réutilisables ;
\textbf{(3)} une variation à jouer. Travaillez une situation par jour pendant 40 jours ---
à voix haute, en jouant les deux rôles.

\begin{astuce}[Le protocole « 5 passes » pour chaque dialogue]
\textbf{Passe 1 :} lisez le dialogue en français seulement --- comprenez la situation.
\textbf{Passe 2 :} lisez-le en anglais à voix haute, lentement, en regardant les
prononciations. \textbf{Passe 3 :} écoutez-le (application de synthèse vocale, ou lisez-le
vous-même en vous enregistrant). \textbf{Passe 4 :} jouez les deux rôles, le livre ouvert.
\textbf{Passe 5 :} jouez-le \textbf{sans le livre}, en changeant trois détails (nom, lieu,
objet). C'est la passe qui compte : elle transforme un texte appris en langage disponible.
\end{astuce}

\section{Situation 1 --- Se présenter (formel)}
\begin{dialoguebox}[1 --- lors d'un rendez-vous professionnel]
\dial{M. Sow}{Good morning. I'm Moussa Sow. Nice to meet you.}{Bonjour. Je suis Moussa Sow. Ravi de vous rencontrer.}
\dial{Mme Clarke}{Good morning, Mr Sow. I'm Helen Clarke. Please, have a seat.}{Bonjour, Monsieur Sow. Je suis Helen Clarke. Asseyez-vous, je vous prie.}
\dial{M. Sow}{Thank you very much.}{Merci beaucoup.}
\dial{Mme Clarke}{So, what do you do, Mr Sow?}{Alors, que faites-vous dans la vie, Monsieur Sow ?}
\dial{M. Sow}{I'm an accountant. I work for a trading company in Dakar.}{Je suis comptable. Je travaille pour une société de négoce à Dakar.}
\dial{Mme Clarke}{How long have you been working there?}{Depuis combien de temps y travaillez-vous ?}
\dial{M. Sow}{For about six years now. I started as an assistant.}{Depuis environ six ans. J'ai commencé comme assistant.}
\dial{Mme Clarke}{And where are you from originally?}{Et d'où venez-vous à l'origine ?}
\dial{M. Sow}{I'm from Thiès, but I've been living in Dakar since 2018.}{Je viens de Thiès, mais j'habite Dakar depuis 2018.}
\dial{Mme Clarke}{Lovely. Shall we start, then?}{Parfait. Commençons-nous alors ?}
\dial{M. Sow}{Of course. I'm ready.}{Bien sûr. Je suis prêt.}
\end{dialoguebox}

\begin{tableau}[Blocs réutilisables]
\begin{tabularx}{\linewidth}{@{}p{6.2cm}X@{}}
\entete{Anglais} & \entete{Français} \tabularnewline
\hline
\en{I'm\ldots\ / My name is\ldots} & Je suis\ldots\ / Je m'appelle\ldots \tabularnewline
\en{Nice to meet you.} & Ravi de vous rencontrer. \tabularnewline
\en{Likewise. / Same here.} & Moi de même. \tabularnewline
\en{I'm a\ldots\ (job). I work for / as / in\ldots} & Je suis\ldots\ Je travaille pour / comme / dans\ldots \tabularnewline
\en{I've been working there for six years.} & J'y travaille depuis six ans. \tabularnewline
\en{I'm from\ldots\ but I live in\ldots} & Je viens de\ldots\ mais j'habite à\ldots \tabularnewline
\en{How do you do?} (très formel) & Comment allez-vous ? (première rencontre) \tabularnewline
\en{Shall we start?} & Commençons-nous ? \tabularnewline
\end{tabularx}
\end{tableau}

\begin{mission}[Variation]
Rejouez le dialogue en vous présentant \textbf{vous}, avec votre vrai métier, votre vraie
ville, votre vraie durée d'expérience. Trois fois de suite, chaque fois plus vite.
\end{mission}

\section{Situation 2 --- Se présenter (informel) et \en{small talk}}
\begin{dialoguebox}[2 --- à une fête de famille]
\dial{Ali}{Hi! I don't think we've met. I'm Ali.}{Salut ! Je ne crois pas qu'on se connaisse. Je suis Ali.}
\dial{Sarah}{Hi Ali, I'm Sarah. How do you know the family?}{Salut Ali, je suis Sarah. Comment connais-tu la famille ?}
\dial{Ali}{Ousmane is my cousin. And you?}{Ousmane est mon cousin. Et toi ?}
\dial{Sarah}{I work with his sister. So, what do you do?}{Je travaille avec sa s\oe ur. Alors, que fais-tu ?}
\dial{Ali}{I'm a teacher --- I teach mathematics at a high school.}{Je suis enseignant --- j'enseigne les mathématiques au lycée.}
\dial{Sarah}{Oh, that must be interesting. How long have you been teaching?}{Oh, ça doit être intéressant. Depuis combien de temps enseignes-tu ?}
\dial{Ali}{About five years. It's hard work, but I love it. And what about you?}{Environ cinq ans. C'est un travail difficile, mais j'adore. Et toi ?}
\dial{Sarah}{I'm a nurse at the regional hospital.}{Je suis infirmière à l'hôpital régional.}
\dial{Ali}{Nice! Do you live around here?}{Super ! Tu habites dans le coin ?}
\dial{Sarah}{Yes, just ten minutes away. By the way, the food is amazing, isn't it?}{Oui, à dix minutes d'ici. Au fait, la nourriture est délicieuse, non ?}
\dial{Ali}{It really is. Have you tried the fish?}{Vraiment. Tu as goûté le poisson ?}
\dial{Sarah}{Not yet --- I will!}{Pas encore --- je vais le faire !}
\end{dialoguebox}

\begin{tableau}[Le \en{small talk} : 12 amorces universelles]
\begin{tabularx}{\linewidth}{@{}p{6.4cm}X@{}}
\entete{Amorce} & \entete{Emploi} \tabularnewline
\hline
\en{Lovely day, isn't it?} & Le temps (sujet sans risque) \tabularnewline
\en{How do you know everyone here?} & Fête, mariage, réunion de famille \tabularnewline
\en{So, what do you do?} & La question standard --- répondez puis \textbf{retournez-la} \tabularnewline
\en{How long have you lived here?} & Résidence \tabularnewline
\en{Did you watch the match last night?} & Sport (le meilleur sujet universel) \tabularnewline
\en{The food is great, isn't it?} & Repas, réception \tabularnewline
\en{Have you been here before?} & Lieu nouveau \tabularnewline
\en{How was your trip?} & Voyage \tabularnewline
\en{What brings you here?} & Raison de la présence (professionnel) \tabularnewline
\en{Do you have any plans for the weekend?} & Vendredi \tabularnewline
\en{How's your week going?} & Lundi ou mardi \tabularnewline
\en{Is this seat taken?} & Transport, café, salle d'attente \tabularnewline
\end{tabularx}
\end{tableau}

\begin{astuce}[La règle des trois relances]
Une conversation meurt quand on répond par une phrase fermée. Entraînez-vous à toujours
ajouter : \textbf{(1)} une réponse, \textbf{(2)} un détail, \textbf{(3)} une question en
retour.\\
\en{``Do you like football?''} --- \en{``Yes, I love it. I play every Saturday with my
brothers. Do you play?''}\\
Trois éléments au lieu d'un mot : la conversation vit.
\end{astuce}

\section{Situation 3 --- Salutations, au revoir, et formules de politesse}
\begin{tableau}[Toutes les salutations, du plus formel au plus familier]
\footnotesize
\begin{longtable}{@{}p{5.2cm}p{4.4cm}p{4.2cm}@{}}
\entete{Rencontrer} & \entete{Prendre congé} & \entete{Contexte} \tabularnewline
\endfirsthead
\entete{Rencontrer} & \entete{Prendre congé} & \entete{Contexte} \tabularnewline
\endhead
\hline
\en{Good morning.} (jusqu'à midi) & \en{Goodbye.} & Formel, universel \tabularnewline
\en{Good afternoon.} (12 h - 18 h) & \en{Have a nice day.} & Formel, commerce \tabularnewline
\en{Good evening.} (après 18 h) & \en{Have a good evening.} & Formel \tabularnewline
\en{Hello. / Hi.} & \en{See you. / See you later.} & Neutre \tabularnewline
\en{Hi there!} & \en{See you tomorrow / on Monday.} & Amical \tabularnewline
\en{Hey!} & \en{Take care.} & Familier, chaleureux \tabularnewline
\en{How are you?} & \en{Catch you later!} & Familier \tabularnewline
\en{How are you doing?} & \en{I'm off --- see you!} & Familier \tabularnewline
\en{How's it going?} & \en{Bye for now.} & Neutre \tabularnewline
\en{What's up? / How come?} & \en{Have a good one.} & Très familier (AmE) \tabularnewline
\en{Long time no see!} & \en{It was great seeing you.} & Retrouvailles \tabularnewline
\en{Pleased to meet you.} & \en{It was a pleasure.} & Première rencontre formelle \tabularnewline
\end{longtable}
\end{tableau}

\begin{tableau}[Répondre à « comment ça va ? »]
\begin{tabularx}{\linewidth}{@{}p{5.4cm}X@{}}
\entete{Réponse} & \entete{Nuance} \tabularnewline
\hline
\en{I'm fine, thanks. And you?} & Standard, poli, toujours correct. \tabularnewline
\en{Pretty good, thanks!} & Positif, naturel. \tabularnewline
\en{Not bad. You?} & Neutre, courant. \tabularnewline
\en{Can't complain.} & « Je ne peux pas me plaindre » --- très idiomatique. \tabularnewline
\en{Same old, same old.} & « Comme d'habitude ». \tabularnewline
\en{I'm doing well, thank you.} & Plus formel. \tabularnewline
\en{Great, thanks! I'm really happy today.} & Enthousiaste. \tabularnewline
\en{To be honest, I'm exhausted.} & Sincère --- ouvre une vraie conversation. \tabularnewline
\en{I've been better, but it's fine.} & Pudeur sur un moment difficile. \tabularnewline
\end{tabularx}
\end{tableau}

\begin{tableau}[Politesse : 20 formules]
\begin{tabularx}{\linewidth}{@{}X X@{}}
\en{Thank you very much.} --- merci beaucoup & \en{Thanks a lot.} --- merci bien \\
\en{I really appreciate it.} --- j'apprécie vraiment & \en{You're welcome.} --- de rien \\
\en{My pleasure.} --- avec plaisir & \en{No problem. / Don't mention it.} --- pas de souci \\
\en{I'm sorry.} --- je suis désolé & \en{Excuse me.} --- pardon (pour attirer l'attention) \\
\en{Sorry to bother you.} --- pardon de vous déranger & \en{Pardon? / Sorry?} --- pardon ? (faire répéter) \\
\en{Could you repeat that, please?} --- pourriez-vous répéter ? & \en{Could you speak more slowly?} --- pourriez-vous parler plus lentement ? \\
\en{Would you mind\ldots?} --- cela vous ennuierait-il de\ldots ? & \en{Could I possibly\ldots?} --- pourrais-je éventuellement\ldots ? \\
\en{I'd appreciate it if you could\ldots} --- je vous serais reconnaissant de\ldots & \en{That's very kind of you.} --- c'est très aimable à vous \\
\en{After you.} --- je vous en prie (passer devant) & \en{Go ahead.} --- allez-y \\
\en{Make yourself at home.} --- faites comme chez vous & \en{Help yourself.} --- servez-vous \\
\en{Congratulations!} --- félicitations ! & \en{Best wishes!} --- meilleurs v\oe ux ! \\
\en{Good luck!} --- bonne chance ! & \en{Take care of yourself.} --- prenez soin de vous \\
\en{Get well soon.} --- bon rétablissement & \en{My condolences.} --- mes condoléances \\
\en{I'm sorry for your loss.} --- je suis désolé pour votre perte & \en{Safe travels!} --- bon voyage ! \\
\end{tabularx}
\end{tableau}

\section{Situation 4 --- Au café et au restaurant}
\begin{dialoguebox}[4 --- commander, demander l'addition]
\dial{Serveur}{Good afternoon! Table for how many?}{Bonjour ! Une table pour combien ?}
\dial{Vous}{Two, please. Do you have a table outside?}{Deux, s'il vous plaît. Avez-vous une table en terrasse ?}
\dial{Serveur}{Yes, right this way. Here are the menus.}{Oui, par ici. Voici les cartes.}
\dial{Vous}{Thank you. What do you recommend?}{Merci. Que recommandez-vous ?}
\dial{Serveur}{The grilled fish is very popular today.}{Le poisson grillé est très demandé aujourd'hui.}
\dial{Vous}{That sounds good. I'll have the grilled fish, please.}{Cela semble bon. Je vais prendre le poisson grillé, s'il vous plaît.}
\dial{Ami}{And I'll take the chicken with rice.}{Et je prendrai le poulet avec du riz.}
\dial{Serveur}{Anything to drink?}{Quelque chose à boire ?}
\dial{Vous}{A large bottle of water and two orange juices, please.}{Une grande bouteille d'eau et deux jus d'orange, s'il vous plaît.}
\dial{Serveur}{Right away. Do you have any allergies?}{Tout de suite. Avez-vous des allergies ?}
\dial{Vous}{No, we're fine, thank you.}{Non, tout va bien, merci.}
\variante{(vingt minutes plus tard)}
\dial{Serveur}{Is everything all right?}{Tout se passe bien ?}
\dial{Vous}{Yes, it's delicious, thank you. Could we have some bread, please?}{Oui, c'est délicieux, merci. Pourrions-nous avoir du pain ?}
\dial{Serveur}{Of course. Enjoy your meal!}{Bien sûr. Bon appétit !}
\dial{Vous}{Could we have the bill, please?}{Pourrions-nous avoir l'addition, s'il vous plaît ?}
\dial{Serveur}{Here you are. Service is not included.}{Voici. Le service n'est pas compris.}
\dial{Vous}{No problem --- we'll leave a tip. Keep the change.}{Pas de problème --- nous laisserons un pourboire. Gardez la monnaie.}
\end{dialoguebox}

\begin{tableau}[Les 15 blocs du restaurant]
\begin{tabularx}{\linewidth}{@{}p{6.4cm}X@{}}
\entete{Bloc} & \entete{Emploi} \tabularnewline
\hline
\en{A table for two, please.} & En arrivant \tabularnewline
\en{Do you have a table by the window?} & Préférence \tabularnewline
\en{I've booked a table under the name\ldots} & Réservation \tabularnewline
\en{Could I see the menu, please?} & Demander la carte \tabularnewline
\en{What's the soup of the day?} & Demander le plat du jour \tabularnewline
\en{What do you recommend?} & Demander un conseil \tabularnewline
\en{I'll have the\ldots, please.} & \textbf{La formule de commande} --- jamais \en{*I want} \tabularnewline
\en{I'd like the\ldots, please.} & Variante polie \tabularnewline
\en{Could I have\ldots\ without onions?} & Exclusion d'un ingrédient \tabularnewline
\en{I'm vegetarian / I don't eat meat.} & Régime \tabularnewline
\en{Could we have some water, please?} & Demande en cours de repas \tabularnewline
\en{It was delicious, thank you.} & Compliment \tabularnewline
\en{Could we have the bill / check, please?} & Addition \tabularnewline
\en{Is service included?} & Question pourboire \tabularnewline
\en{Can I pay by card? / Keep the change.} & Paiement \tabularnewline
\end{tabularx}
\end{tableau}

\begin{piege}[\en{*I want a coffee}]
Grammaticalement correct, socialement brutal. En anglais, \en{I want} sonne comme un ordre
d'enfant. Utilisez \en{I'd like\ldots}, \en{I'll have\ldots}, \en{Could I have\ldots?}
--- trois secondes de politesse qui changent tout le ton d'un échange.
\end{piege}

\section{Situation 5 --- Au marché et dans les magasins}
\begin{dialoguebox}[5 --- acheter, demander le prix, négocier]
\dial{Vendeur}{Hello! Can I help you find something?}{Bonjour ! Je peux vous aider à trouver quelque chose ?}
\dial{Vous}{Yes, please. I'm looking for a shirt --- a blue one.}{Oui, s'il vous plaît. Je cherche une chemise --- une bleue.}
\dial{Vendeur}{We have this one, and this one in a lighter blue.}{Nous avons celle-ci, et celle-là en bleu plus clair.}
\dial{Vous}{How much is this one?}{Combien coûte celle-ci ?}
\dial{Vendeur}{It's 12,000 francs.}{C'est 12\,000 francs.}
\dial{Vous}{That's a bit expensive. Can you do 10,000?}{C'est un peu cher. Pouvez-vous la faire à 10\,000 ?}
\dial{Vendeur}{Tell you what --- 11,000, and it's a very good price.}{Voici ce que je vous propose --- 11\,000, et c'est un très bon prix.}
\dial{Vous}{Deal. Can I try it on first?}{D'accord. Puis-je l'essayer d'abord ?}
\dial{Vendeur}{Of course, the fitting room is over there.}{Bien sûr, la cabine d'essayage est par là.}
\dial{Vous}{It fits perfectly. I'll take it.}{Elle me va parfaitement. Je la prends.}
\dial{Vendeur}{Would you like a bag? Do you need a receipt?}{Voulez-vous un sac ? Avez-vous besoin d'un reçu ?}
\dial{Vous}{A bag, please. No receipt --- I'm paying in cash.}{Un sac, s'il vous plaît. Pas de reçu --- je paie en espèces.}
\end{dialoguebox}

\begin{tableau}[Achats : 18 blocs]
\begin{tabularx}{\linewidth}{@{}p{6.4cm}X@{}}
\entete{Bloc} & \entete{Emploi} \tabularnewline
\hline
\en{I'm just looking, thanks.} & Refuser poliment de l'aide \tabularnewline
\en{I'm looking for\ldots} & Annoncer sa recherche \tabularnewline
\en{Do you have this in a bigger / smaller size?} & Tailles \tabularnewline
\en{Do you have it in another colour?} & Couleur \tabularnewline
\en{Where can I find\ldots?} & Localisation dans le magasin \tabularnewline
\en{How much is it? / How much does it cost?} & Prix \tabularnewline
\en{That's too expensive for me.} & Négocier \tabularnewline
\en{Can you give me a discount?} & Demander une remise \tabularnewline
\en{What's your best price?} & Négocier \tabularnewline
\en{Can I try it on?} & Essayer \tabularnewline
\en{Where is the fitting room?} & Cabine \tabularnewline
\en{It fits well / It's too tight / too loose.} & Essayage \tabularnewline
\en{It suits you!} & Compliment (aller à la personne) \tabularnewline
\en{I'll take it.} & Acheter \tabularnewline
\en{Can I pay by card? / Is cash okay?} & Paiement \tabularnewline
\en{Could I have a receipt, please?} & Reçu \tabularnewline
\en{Can I get a refund / exchange it?} & Retour, échange \tabularnewline
\en{Is it still under guarantee / warranty?} & Garantie \tabularnewline
\end{tabularx}
\end{tableau}

\section{Situation 6 --- Demander son chemin}
\begin{dialoguebox}[6 --- dans la rue]
\dial{Vous}{Excuse me, sorry to bother you. How do I get to the central market?}{Pardon, excusez-moi de vous déranger. Comment aller au marché central ?}
\dial{Passant}{The central market? It's about a fifteen-minute walk from here.}{Le marché central ? C'est à environ quinze minutes à pied d'ici.}
\dial{Vous}{Oh, I see. Which way should I go?}{Ah, d'accord. Par où dois-je aller ?}
\dial{Passant}{Go straight on for two blocks, then turn left at the traffic lights.}{Allez tout droit pendant deux rues, puis tournez à gauche au feu.}
\dial{Vous}{Straight on, then left at the lights. Is it far from the mosque?}{Tout droit, puis à gauche au feu. C'est loin de la mosquée ?}
\dial{Passant}{No, it's just opposite the mosque, next to the pharmacy.}{Non, c'est juste en face de la mosquée, à côté de la pharmacie.}
\dial{Vous}{Great. And is there a bus that goes there?}{Parfait. Et y a-t-il un bus qui y va ?}
\dial{Passant}{Yes, bus number 4 stops right in front of it.}{Oui, le bus numéro 4 s'arrête juste devant.}
\dial{Vous}{Thank you very much for your help.}{Merci beaucoup pour votre aide.}
\dial{Passant}{No problem. You can't miss it.}{Pas de souci. Vous ne pouvez pas le rater.}
\end{dialoguebox}

\begin{tableau}[Les indications : tout le vocabulaire du déplacement]
\begin{tabularx}{\linewidth}{@{}p{5.4cm}X@{}}
\entete{Instruction} & \entete{Sens} \tabularnewline
\hline
\en{Go straight on / ahead.} & Allez tout droit. \tabularnewline
\en{Turn left / right.} & Tournez à gauche / à droite. \tabularnewline
\en{Take the first / second turning.} & Prenez la première / deuxième rue. \tabularnewline
\en{At the roundabout, take the second exit.} & Au rond-point, prenez la deuxième sortie. \tabularnewline
\en{Cross the street / the bridge.} & Traversez la rue / le pont. \tabularnewline
\en{Go past the bank.} & Passez devant la banque. \tabularnewline
\en{It's on your left / right.} & C'est sur votre gauche / droite. \tabularnewline
\en{It's next to / opposite / behind / in front of\ldots} & C'est à côté de / en face de / derrière / devant\ldots \tabularnewline
\en{It's at the end of the street.} & C'est au bout de la rue. \tabularnewline
\en{It's about five minutes' walk.} & C'est à cinq minutes à pied. \tabularnewline
\en{You can't miss it.} & Vous ne pouvez pas le rater. \tabularnewline
\en{I'm afraid I'm not from around here.} & Je ne suis pas du coin. \tabularnewline
\en{Let me check on my phone.} & Laissez-moi vérifier sur mon téléphone. \tabularnewline
\end{tabularx}
\end{tableau}

\section{Situation 7 --- À l'hôtel}
\begin{dialoguebox}[7 --- arrivée et problème de chambre]
\dial{Réceptionniste}{Good evening, sir. How can I help you?}{Bonsoir, monsieur. Comment puis-je vous aider ?}
\dial{Vous}{Good evening. I have a reservation under the name Diop, for three nights.}{Bonsoir. J'ai une réservation au nom de Diop, pour trois nuits.}
\dial{Réceptionniste}{Let me check\ldots\ Yes, here it is. A single room with breakfast. Could I see your ID, please?}{Je vérifie\ldots\ Oui, la voici. Une chambre simple avec petit-déjeuner. Puis-je voir votre pièce d'identité ?}
\dial{Vous}{Of course, here's my passport.}{Bien sûr, voici mon passeport.}
\dial{Réceptionniste}{Thank you. You're in room 304, on the third floor. Here's your key. Breakfast is from 7 to 10 am.}{Merci. Vous êtes en chambre 304, au troisième étage. Voici votre clé. Le petit-déjeuner est servi de 7 h à 10 h.}
\dial{Vous}{Is there Wi-Fi in the rooms?}{Y a-t-il du wi-fi dans les chambres ?}
\dial{Réceptionniste}{Yes, it's free. The password is on your key card holder.}{Oui, il est gratuit. Le mot de passe est sur l'étui de votre carte.}
\dial{Vous}{Perfect, thank you.}{Parfait, merci.}
\variante{(le lendemain matin, à la réception)}
\dial{Vous}{Excuse me, I'm sorry, but there's a problem with my room.}{Excusez-moi, je suis désolé, mais il y a un problème avec ma chambre.}
\dial{Réceptionniste}{I'm sorry to hear that. What seems to be the problem?}{Je suis navrée de l'apprendre. Quel est le problème ?}
\dial{Vous}{The air conditioning isn't working and there's no hot water.}{La climatisation ne fonctionne pas et il n'y a pas d'eau chaude.}
\dial{Réceptionniste}{I do apologise. I'll send someone up right away. If we can't fix it, we'll move you to another room.}{Je vous présente mes excuses. J'envoie quelqu'un tout de suite. Si nous ne pouvons pas réparer, nous vous changerons de chambre.}
\dial{Vous}{That would be great. Thank you for your help.}{Ce serait parfait. Merci pour votre aide.}
\end{dialoguebox}

\begin{tableau}[Se plaindre sans agressivité --- la technique du sandwich]
\begin{tabularx}{\linewidth}{@{}p{6.2cm}X@{}}
\entete{Étape} & \entete{Formule} \tabularnewline
\hline
1. Excuse polie & \en{Excuse me, I'm sorry, but\ldots} / \en{I'm afraid there's a problem with\ldots} \tabularnewline
2. Fait précis & \en{The air conditioning isn't working.} / \en{I've been waiting for 40 minutes.} / \en{This isn't what I ordered.} \tabularnewline
3. Demande claire & \en{Could you send someone, please?} / \en{Is there anything you can do?} / \en{Could I have a refund?} \tabularnewline
4. Remerciement & \en{Thank you, I really appreciate it.} \tabularnewline
\end{tabularx}
\end{tableau}

\begin{astuce}[\en{I'm afraid} et \en{I'm sorry} adoucissent tout]
\en{I'm afraid there's a mistake.} (Je crains qu'il y ait une erreur.) sonne bien mieux que
\en{There's a mistake.} En anglais, la plainte commence presque toujours par une excuse ---
même quand c'est vous qui avez raison. C'est culturel, pas logique.
\end{astuce}

\section{Situation 8 --- À l'aéroport}
\begin{dialoguebox}[8 --- enregistrement et contrôle]
\dial{Agent}{Good morning. Where are you flying to today?}{Bonjour. Où vous envolez-vous aujourd'hui ?}
\dial{Vous}{To London, via Casablanca.}{Vers Londres, via Casablanca.}
\dial{Agent}{May I see your passport and ticket, please? Do you have any luggage to check in?}{Puis-je voir votre passeport et votre billet ? Avez-vous des bagages à enregistrer ?}
\dial{Vous}{Here you are. Just one suitcase --- this is my hand luggage.}{Voici. Une seule valise --- ceci est mon bagage à main.}
\dial{Agent}{Thank you. Do you have anything sharp or liquid in your hand luggage?}{Merci. Avez-vous des objets tranchants ou des liquides dans votre bagage à main ?}
\dial{Vous}{No, nothing like that.}{Non, rien de ce genre.}
\dial{Agent}{Great. Window or aisle seat?}{Parfait. Siège côté hublot ou côté couloir ?}
\dial{Vous}{A window seat, please, if possible.}{Un siège côté hublot, s'il vous plaît, si c'est possible.}
\dial{Agent}{You're on seat 24A. Boarding starts at gate 12 at 10:40.}{Vous êtes au siège 24A. L'embarquement commence à la porte 12 à 10 h 40.}
\dial{Vous}{Thank you. How long is the flight?}{Merci. Combien de temps dure le vol ?}
\dial{Agent}{About six hours. Have a pleasant flight!}{Environ six heures. Bon vol !}
\variante{(au contrôle des passeports)}
\dial{Officier}{Passport, please. What's the purpose of your visit?}{Votre passeport, s'il vous plaît. Quel est le motif de votre visite ?}
\dial{Vous}{It's for a training course. I'll be here for two weeks.}{C'est pour une formation. Je resterai deux semaines.}
\dial{Officier}{Where will you be staying?}{Où logerez-vous ?}
\dial{Vous}{At the Hilton Hotel, in the city centre.}{À l'hôtel Hilton, au centre-ville.}
\dial{Officier}{Thank you. Enjoy your stay.}{Merci. Bon séjour.}
\end{dialoguebox}

\section{Situation 9 --- Chez le médecin}
\begin{dialoguebox}[9 --- consultation]
\dial{Médecin}{Good morning. What seems to be the problem?}{Bonjour. Quel est le problème ?}
\dial{Vous}{I've been feeling ill for three days. I have a headache and a sore throat.}{Je me sens malade depuis trois jours. J'ai mal à la tête et mal à la gorge.}
\dial{Médecin}{Do you have a fever?}{Avez-vous de la fièvre ?}
\dial{Vous}{Yes, I think so. I felt hot last night.}{Oui, je crois. J'avais chaud hier soir.}
\dial{Médecin}{Let me take your temperature\ldots\ 38.5. Open your mouth and say ``ah''. Any cough?}{Je prends votre température\ldots\ 38,5. Ouvrez la bouche et dites « ah ». De la toux ?}
\dial{Vous}{A dry cough, mostly at night.}{Une toux sèche, surtout la nuit.}
\dial{Médecin}{Are you allergic to any medication? Do you take any other medicines?}{Êtes-vous allergique à un médicament ? Prenez-vous d'autres traitements ?}
\dial{Vous}{No allergies, and I only take vitamins.}{Aucune allergie, et je ne prends que des vitamines.}
\dial{Médecin}{It looks like a bad cold, nothing serious. I'll prescribe something for the fever and the cough. Rest for three days and drink plenty of water.}{On dirait un gros rhume, rien de grave. Je vais vous prescrire quelque chose pour la fièvre et la toux. Reposez-vous trois jours et buvez beaucoup d'eau.}
\dial{Vous}{Thank you, doctor. Should I come back if it doesn't get better?}{Merci, docteur. Dois-je revenir si cela ne s'améliore pas ?}
\dial{Médecin}{Yes, come back in a week if you're not better. Take care.}{Oui, revenez dans une semaine si vous n'allez pas mieux. Prenez soin de vous.}
\end{dialoguebox}

\begin{tableau}[Décrire ses symptômes --- 20 blocs]
\begin{tabularx}{\linewidth}{@{}p{6.4cm}X@{}}
\entete{Bloc} & \entete{Sens} \tabularnewline
\hline
\en{I don't feel well.} & Je ne me sens pas bien. \tabularnewline
\en{I've been feeling ill for three days.} & Je suis malade depuis trois jours. \tabularnewline
\en{I have a headache / stomachache / backache.} & J'ai mal à la tête / au ventre / au dos. \tabularnewline
\en{I have a sore throat.} & J'ai mal à la gorge. \tabularnewline
\en{I have a runny nose and a cough.} & J'ai le nez qui coule et je tousse. \tabularnewline
\en{I have a fever / a temperature.} & J'ai de la fièvre. \tabularnewline
\en{I feel dizzy / sick.} & J'ai la tête qui tourne / je me sens mal. \tabularnewline
\en{It hurts when I swallow.} & Ça fait mal quand j'avale. \tabularnewline
\en{My leg is swollen.} & Ma jambe est enflée. \tabularnewline
\en{I've been vomiting.} & J'ai vomi. \tabularnewline
\en{I can't sleep.} & Je n'arrive pas à dormir. \tabularnewline
\en{I'm always tired.} & Je suis toujours fatigué. \tabularnewline
\en{I have an allergy to\ldots} & Je suis allergique à\ldots \tabularnewline
\en{I take medicine for\ldots} & Je prends un traitement pour\ldots \tabularnewline
\en{I need a prescription.} & J'ai besoin d'une ordonnance. \tabularnewline
\en{I need a medical certificate.} & J'ai besoin d'un certificat médical. \tabularnewline
\en{How long will it take to recover?} & Combien de temps pour guérir ? \tabularnewline
\en{Is it serious?} & Est-ce grave ? \tabularnewline
\en{Should I see a specialist?} & Dois-je voir un spécialiste ? \tabularnewline
\en{What should I avoid?} & Que dois-je éviter ? \tabularnewline
\end{tabularx}
\end{tableau}

\section{Situation 10 --- Au téléphone}
\begin{dialoguebox}[10 --- appel professionnel]
\dial{Vous}{Good morning, Diop and Sons, how can I help you?}{Bonjour, Diop et Fils, comment puis-je vous aider ?}
\dial{Appelant}{Hello, this is John Mensah from Bright Ltd. May I speak to Mr Diop, please?}{Bonjour, ici John Mensah de Bright Ltd. Puis-je parler à Monsieur Diop ?}
\dial{Vous}{I'm afraid he's in a meeting at the moment. Can I take a message?}{Je crains qu'il ne soit en réunion pour le moment. Puis-je prendre un message ?}
\dial{Appelant}{Yes, please. Could you ask him to call me back about the delivery?}{Oui, s'il vous plaît. Pourriez-vous lui demander de me rappeler au sujet de la livraison ?}
\dial{Vous}{Of course. Could you spell your surname for me, please?}{Bien sûr. Pourriez-vous épeler votre nom de famille ?}
\dial{Appelant}{It's M-E-N-S-A-H.}{C'est M-E-N-S-A-H.}
\dial{Vous}{Thank you. And what's the best number to reach you on?}{Merci. Et quel est le meilleur numéro pour vous joindre ?}
\dial{Appelant}{It's 77 123 45 67. I'm available after 2 pm.}{C'est le 77 123 45 67. Je suis disponible après 14 h.}
\dial{Vous}{Right, I'll make sure he gets the message. Anything else?}{D'accord, je m'assure qu'il reçoive le message. Autre chose ?}
\dial{Appelant}{No, that's all. Thanks for your help.}{Non, c'est tout. Merci de votre aide.}
\dial{Vous}{You're welcome. Goodbye, Mr Mensah.}{De rien. Au revoir, Monsieur Mensah.}
\end{dialoguebox}

\begin{tableau}[Le téléphone : 20 blocs]
\begin{tabularx}{\linewidth}{@{}p{6.4cm}X@{}}
\entete{Bloc} & \entete{Emploi} \tabularnewline
\hline
\en{Hello, this is Moussa speaking.} & Se présenter au téléphone (\textbf{jamais} \en{*I am Moussa}) \tabularnewline
\en{May I speak to\ldots, please?} & Demander quelqu'un \tabularnewline
\en{Could I ask who's calling?} & Demander qui appelle \tabularnewline
\en{Hold on, please. / One moment, please.} & Faire patienter \tabularnewline
\en{I'm afraid he's not available.} & Personne n'est joignable \tabularnewline
\en{Can I take a message?} & Prendre un message \tabularnewline
\en{Could you ask him to call me back?} & Demander un rappel \tabularnewline
\en{I'll call back later.} & Je rappellerai plus tard. \tabularnewline
\en{Could you spell that for me?} & Faire épeler \tabularnewline
\en{Sorry, you're breaking up.} & Vous grésillez / ça coupe. \tabularnewline
\en{The line is very bad. Can you hear me?} & La ligne est mauvaise. \tabularnewline
\en{I'm sorry, could you repeat that?} & Faire répéter \tabularnewline
\en{Could you speak a bit louder?} & Plus fort \tabularnewline
\en{Let me put you through.} & Je vous passe quelqu'un. \tabularnewline
\en{I think we've been cut off.} & Nous avons été coupés. \tabularnewline
\en{I'll send you an email to confirm.} & Je confirme par écrit. \tabularnewline
\en{Thanks for calling. Goodbye.} & Conclure \tabularnewline
\en{Is now a good time to talk?} & Vérifier la disponibilité \tabularnewline
\en{I'm calling about\ldots} & Annoncer l'objet de l'appel \tabularnewline
\en{Sorry to call so late.} & S'excuser de l'heure \tabularnewline
\end{tabularx}
\end{tableau}

\section{Situation 11 --- L'entretien d'embauche}
\begin{dialoguebox}[11 --- les six questions incontournables]
\dial{Recruteur}{Thanks for coming in. Could you tell me a little about yourself?}{Merci d'être venu. Pouvez-vous vous présenter brièvement ?}
\dial{Vous}{Of course. I'm an accountant with six years' experience. I've been working for a trading company in Dakar since 2020, where I handle invoicing and monthly reports.}{Bien sûr. Je suis comptable avec six ans d'expérience. Je travaille pour une société de négoce à Dakar depuis 2020, où je m'occupe de la facturation et des rapports mensuels.}
\dial{Recruteur}{Why do you want to work for us?}{Pourquoi souhaitez-vous travailler chez nous ?}
\dial{Vous}{Your company is growing fast in West Africa, and I'd like to work on larger accounts. I also value your training programmes.}{Votre entreprise se développe rapidement en Afrique de l'Ouest, et j'aimerais travailler sur des comptes plus importants. J'apprécie aussi vos programmes de formation.}
\dial{Recruteur}{What are your strengths and weaknesses?}{Quelles sont vos forces et vos faiblesses ?}
\dial{Vous}{My strength is accuracy --- I rarely make errors in reporting. My weakness used to be public speaking, but I've been taking courses, and I've improved a lot.}{Ma force est la rigueur --- je fais rarement des erreurs dans les rapports. Ma faiblesse était la prise de parole en public, mais je suis des formations et je me suis beaucoup amélioré.}
\dial{Recruteur}{Where do you see yourself in five years?}{Où vous voyez-vous dans cinq ans ?}
\dial{Vous}{Ideally, leading a small team and helping the company expand into new markets.}{Idéalement, à la tête d'une petite équipe, en aidant l'entreprise à s'étendre vers de nouveaux marchés.}
\dial{Recruteur}{Why did you leave your last job?}{Pourquoi avez-vous quitté votre dernier emploi ?}
\dial{Vous}{There were no opportunities for promotion, so I'm looking for a new challenge. I left on very good terms.}{Il n'y avait pas de perspectives de promotion, je cherche donc un nouveau défi. Je suis parti en très bons termes.}
\dial{Recruteur}{Do you have any questions for us?}{Avez-vous des questions à nous poser ?}
\dial{Vous}{Yes --- what does a typical week look like in this role, and what would success look like after six months?}{Oui --- à quoi ressemble une semaine type dans ce poste, et à quoi ressemblerait la réussite après six mois ?}
\end{dialoguebox}

\begin{astuce}[Les trois réponses à préparer par écrit]
Avant tout entretien, rédigez et mémorisez \textbf{60 secondes} sur : (1) qui vous êtes
professionnellement ; (2) pourquoi cette entreprise ; (3) votre plus belle réussite, avec un
chiffre. Ces trois blocs couvrent 60\,\% de l'entretien. Entraînez-les à voix haute dix fois
--- et enregistrez-vous.
\end{astuce}

\section{Situation 12 --- Donner son opinion et débattre}
\begin{dialoguebox}[12 --- discussion entre amis]
\dial{Ami}{What do you think about working abroad? Everyone wants to leave.}{Que penses-tu de travailler à l'étranger ? Tout le monde veut partir.}
\dial{Vous}{Well, it depends. On the one hand, it's a great opportunity to learn and earn more.}{Eh bien, cela dépend. D'un côté, c'est une grande occasion d'apprendre et de gagner plus.}
\dial{Ami}{Exactly! And you can support your family back home.}{Exactement ! Et tu peux aider ta famille restée au pays.}
\dial{Vous}{That's true. But on the other hand, I think we also need skilled people here. If everyone leaves, who builds the country?}{C'est vrai. Mais d'un autre côté, je pense que nous avons aussi besoin de gens qualifiés ici. Si tout le monde part, qui construit le pays ?}
\dial{Ami}{I see your point, but honestly, without opportunities here, people have no choice.}{Je comprends ton point de vue, mais honnêtement, sans opportunités ici, les gens n'ont pas le choix.}
\dial{Vous}{I couldn't agree more on that. That's why I think the real question is how to create opportunities here.}{Sur ce point, je suis tout à fait d'accord. C'est pourquoi je pense que la vraie question est de savoir comment créer des opportunités ici.}
\dial{Ami}{Fair enough. We're not going to solve that today, though!}{C'est juste. Mais nous n'allons pas résoudre ça aujourd'hui !}
\dial{Vous}{No, we're not --- but it was a good conversation.}{Non, en effet --- mais c'était une bonne conversation.}
\end{dialoguebox}

\begin{tableau}[Exprimer, nuancer, contredire --- 25 blocs]
\begin{tabularx}{\linewidth}{@{}p{6.4cm}X@{}}
\entete{Bloc} & \entete{Fonction} \tabularnewline
\hline
\en{In my opinion, \ldots} & donner un avis neutre \tabularnewline
\en{I think / I believe / I feel that\ldots} & avis courant \tabularnewline
\en{As far as I'm concerned, \ldots} & avis personnel marqué \tabularnewline
\en{From my point of view, \ldots} & point de vue \tabularnewline
\en{I'd say that\ldots} & je dirais que \tabularnewline
\en{It seems to me that\ldots} & il me semble que \tabularnewline
\en{I'm convinced that\ldots} & je suis convaincu que \tabularnewline
\en{I couldn't agree more.} & je suis totalement d'accord \tabularnewline
\en{That's exactly what I think.} & c'est exactement ce que je pense \tabularnewline
\en{You've got a point there.} & tu marques un point \tabularnewline
\en{Fair enough.} & c'est juste (accord partiel) \tabularnewline
\en{I agree up to a point, but\ldots} & je suis d'accord jusqu'à un certain point \tabularnewline
\en{I see what you mean, but\ldots} & je vois ce que tu veux dire, mais \tabularnewline
\en{That's not quite how I see it.} & ce n'est pas tout à fait ainsi que je vois les choses \tabularnewline
\en{I'm afraid I have to disagree.} & je crains de devoir être en désaccord \tabularnewline
\en{With all due respect, \ldots} & avec tout le respect que je vous dois \tabularnewline
\en{On the one hand\ldots on the other hand\ldots} & d'un côté\ldots de l'autre \tabularnewline
\en{It depends on the situation.} & cela dépend de la situation \tabularnewline
\en{That's a good question --- let me think.} & bonne question --- laisse-moi réfléchir \tabularnewline
\en{I'm not sure, but I'd guess that\ldots} & je ne suis pas sûr, mais je suppose que \tabularnewline
\en{To be honest, \ldots} & pour être honnête \tabularnewline
\en{In other words, \ldots} & autrement dit \tabularnewline
\en{Let's agree to disagree.} & restons sur nos positions \tabularnewline
\en{That reminds me of something.} & cela me fait penser à quelque chose \tabularnewline
\en{Going back to what you said earlier\ldots} & pour revenir à ce que tu disais \tabularnewline
\end{tabularx}
\end{tableau}

\section{Situation 13 --- Raconter une histoire}
\begin{dialoguebox}[13 --- une anecdote (passé simple, continu et perfect)]
\dial{Ami}{Have you ever had a really embarrassing moment?}{T'est-il déjà arrivé un moment vraiment embarrassant ?}
\dial{Vous}{Oh yes! Last year, I was travelling to a conference in Banjul.}{Oh oui ! L'année dernière, je voyageais pour une conférence à Banjul.}
\dial{Ami}{What happened?}{Que s'est-il passé ?}
\dial{Vous}{Well, I had arrived early, so I was waiting at the gate. Suddenly, a man started waving at me and shouting my name.}{Eh bien, j'étais arrivé tôt, donc j'attendais à la porte. Soudain, un homme a commencé à me faire signe et à crier mon nom.}
\dial{Ami}{Did you know him?}{Le connaissais-tu ?}
\dial{Vous}{Not at all! It turned out he was waving at his brother, who was standing right behind me.}{Pas du tout ! Il s'est avéré qu'il faisait signe à son frère, qui se tenait juste derrière moi.}
\dial{Ami}{Ha! What did you do?}{Ha ! Qu'as-tu fait ?}
\dial{Vous}{I waved back for about ten seconds before I realised. I've never been so embarrassed in my life.}{J'ai répondu en faisant signe pendant environ dix secondes avant de réaliser. Je n'ai jamais été aussi embarrassé de ma vie.}
\dial{Ami}{That's a great story. Did you ever meet him again?}{C'est une bonne histoire. L'as-tu revu ?}
\dial{Vous}{Actually, yes --- we were on the same plane, and we ended up sitting next to each other!}{En fait oui --- nous étions dans le même avion, et nous avons fini par nous asseoir l'un à côté de l'autre !}
\end{dialoguebox}

\begin{astuce}[La structure d'une bonne anecdote --- 5 étapes]
\textbf{1. Cadre} (passé continu / \en{past perfect}) : \en{It was raining and I was
waiting for the bus.}\\
\textbf{2. Élément déclencheur} (\en{suddenly}, passé simple) : \en{Suddenly, a dog ran
towards me.}\\
\textbf{3. Réaction} : \en{I didn't know what to do, so I climbed onto a bench.}\\
\textbf{4. Résultat} : \en{In the end, the owner came and took it away.}\\
\textbf{5. Commentaire} (\en{present perfect}) : \en{I've never been so scared!}\\
Apprenez cette structure et vous pourrez raconter n'importe quelle histoire de votre vie
--- c'est exactement ce que font les natifs.
\end{astuce}

\section{Situation 14 --- Parler de son pays et de sa culture}
\begin{dialoguebox}[14 --- avec un étranger curieux]
\dial{Étranger}{So you're from Senegal --- what's it like?}{Tu viens donc du Sénégal --- comment c'est ?}
\dial{Vous}{It's a beautiful country on the west coast of Africa. The people are very welcoming, and we have a long tradition of hospitality --- we call it \en{teranga}.}{C'est un beau pays sur la côte ouest de l'Afrique. Les gens sont très accueillants, et nous avons une longue tradition d'hospitalité --- nous l'appelons la teranga.}
\dial{Étranger}{What's the food like? I've never tried Senegalese food.}{Comment est la nourriture ? Je n'ai jamais goûté la cuisine sénégalaise.}
\dial{Vous}{You should try \en{thieboudienne} --- it's fish with rice and vegetables. It's our national dish. We also eat a lot of millet, peanuts and fresh fish.}{Tu devrais essayer le thieboudienne --- c'est du poisson avec du riz et des légumes. C'est notre plat national. Nous mangeons aussi beaucoup de mil, d'arachide et de poisson frais.}
\dial{Étranger}{That sounds delicious. And what languages do people speak?}{Cela semble délicieux. Et quelles langues les gens parlent-ils ?}
\dial{Vous}{French is the official language, but most people speak Wolof at home. There are also Pulaar, Serer, Diola and others.}{Le français est la langue officielle, mais la plupart des gens parlent wolof à la maison. Il y a aussi le pulaar, le sérère, le diola et d'autres.}
\dial{Étranger}{Wow. Is it hot all year round?}{Ouah. Il fait chaud toute l'année ?}
\dial{Vous}{Pretty much. We have a dry season from November to May and a rainy season from June to October. It can get very hot in April.}{À peu près. Nous avons une saison sèche de novembre à mai et une saison des pluies de juin à octobre. Il peut faire très chaud en avril.}
\dial{Étranger}{I'd love to visit one day.}{J'aimerais y aller un jour.}
\dial{Vous}{You'd be very welcome. I'd be happy to show you around.}{Tu serais le bienvenu. Je serais heureux de te faire visiter.}
\end{dialoguebox}

\begin{tableau}[Parler de son pays : 20 blocs]
\begin{tabularx}{\linewidth}{@{}p{6.6cm}X@{}}
\entete{Bloc} & \entete{Sens} \tabularnewline
\hline
\en{I'm from\ldots, in West Africa.} & Je viens de\ldots, en Afrique de l'Ouest. \tabularnewline
\en{It's a country of about 17 million people.} & C'est un pays d'environ 17 millions d'habitants. \tabularnewline
\en{The capital is\ldots, but I'm from\ldots} & La capitale est\ldots, mais je viens de\ldots \tabularnewline
\en{We speak\ldots\ at home and\ldots\ at school.} & Nous parlons\ldots à la maison et\ldots à l'école. \tabularnewline
\en{Our national dish is called\ldots} & Notre plat national s'appelle\ldots \tabularnewline
\en{It's made with rice, fish and vegetables.} & C'est fait avec du riz, du poisson et des légumes. \tabularnewline
\en{The climate is hot and dry for most of the year.} & Le climat est chaud et sec la majeure partie de l'année. \tabularnewline
\en{We have a dry season and a rainy season.} & Nous avons une saison sèche et une saison des pluies. \tabularnewline
\en{Football is by far the most popular sport.} & Le football est de loin le sport le plus populaire. \tabularnewline
\en{Traditional wrestling is also very popular.} & La lutte traditionnelle est aussi très populaire. \tabularnewline
\en{Family is at the centre of our culture.} & La famille est au c\oe ur de notre culture. \tabularnewline
\en{People are known for their hospitality.} & Les gens sont connus pour leur hospitalité. \tabularnewline
\en{Most people are Muslim, and there are also many Christians.} & La plupart des gens sont musulmans, et il y a aussi beaucoup de chrétiens. \tabularnewline
\en{We celebrate\ldots, which is a very important festival.} & Nous célébrons\ldots, qui est une fête très importante. \tabularnewline
\en{The economy is based on agriculture, fishing and services.} & L'économie repose sur l'agriculture, la pêche et les services. \tabularnewline
\en{Tourism is growing, especially around\ldots} & Le tourisme se développe, surtout autour de\ldots \tabularnewline
\en{It's about two hours from the capital by road.} & C'est à environ deux heures de la capitale par la route. \tabularnewline
\en{If you visit, you should definitely see\ldots} & Si tu visites, tu dois absolument voir\ldots \tabularnewline
\en{The best time to come is between November and February.} & La meilleure période pour venir est entre novembre et février. \tabularnewline
\en{I'd be happy to show you around.} & Je serais heureux de te faire visiter. \tabularnewline
\end{tabularx}
\end{tableau}

\section{Situation 15 --- La réunion de travail}
\begin{dialoguebox}[15 --- animer et participer]
\dial{Animatrice}{Good morning, everyone. Thanks for coming. The purpose of today's meeting is to review the second-quarter results. Let's get started.}{Bonjour à tous. Merci d'être venus. L'objet de la réunion d'aujourd'hui est de passer en revue les résultats du deuxième trimestre. Commençons.}
\dial{Vous}{Before we start, could I just clarify the agenda? Are we discussing the budget today as well?}{Avant de commencer, puis-je préciser l'ordre du jour ? Discutons-nous aussi du budget aujourd'hui ?}
\dial{Animatrice}{Good question. We'll come to the budget at the end, if we have time. First, Moussa, could you give us an overview of sales?}{Bonne question. Nous aborderons le budget à la fin, si nous avons le temps. D'abord, Moussa, pourriez-vous nous donner un aperçu des ventes ?}
\dial{Vous}{Sure. Sales increased by 12\,\% compared to last year, which is above our target. However, costs also rose, so the margin is slightly lower.}{Bien sûr. Les ventes ont augmenté de 12\,\% par rapport à l'année dernière, ce qui dépasse notre objectif. Cependant, les coûts ont aussi augmenté, donc la marge est légèrement plus basse.}
\dial{Collègue}{Sorry to interrupt, but does that include the new market?}{Pardon de vous interrompre, mais cela inclut-il le nouveau marché ?}
\dial{Vous}{No, it doesn't. If we include it, the increase is closer to 18\,\%.}{Non. Si on l'inclut, l'augmentation est plus proche de 18\,\%.}
\dial{Animatrice}{That's very encouraging. Let's move on to the next point.}{C'est très encourageant. Passons au point suivant.}
\dial{Collègue}{Could you go back to the costs for a moment? I'd like to understand where they come from.}{Pourriez-vous revenir sur les coûts un instant ? J'aimerais comprendre d'où ils viennent.}
\dial{Vous}{Of course. I'll send you the detailed breakdown after the meeting.}{Bien sûr. Je vous enverrai le détail après la réunion.}
\dial{Animatrice}{Right, to sum up: sales are up, costs need attention, and Moussa will send the details. Any other business?}{Bien, pour résumer : les ventes sont en hausse, les coûts demandent attention, et Moussa enverra le détail. Autres questions ?}
\dial{Vous}{Just one thing --- could we schedule a follow-up meeting next week?}{Une seule chose --- pourrions-nous prévoir une réunion de suivi la semaine prochaine ?}
\dial{Animatrice}{Good idea. I'll send an invitation. Thanks everyone --- meeting adjourned.}{Bonne idée. J'enverrai une invitation. Merci à tous --- la séance est levée.}
\end{dialoguebox}

\begin{tableau}[Animer et participer : 25 blocs]
\begin{tabularx}{\linewidth}{@{}p{6.6cm}X@{}}
\entete{Bloc} & \entete{Fonction} \tabularnewline
\hline
\en{Let's get started.} & Ouvrir \tabularnewline
\en{The purpose of this meeting is to\ldots} & Annoncer l'objet \tabularnewline
\en{There are three items on the agenda.} & Annoncer le plan \tabularnewline
\en{First of all, I'd like to\ldots} & Premier point \tabularnewline
\en{Let's move on to the next point.} & Passer au suivant \tabularnewline
\en{Could you give us an overview of\ldots?} & Demander un aperçu \tabularnewline
\en{To put it simply, \ldots} & Simplifier \tabularnewline
\en{The figures show that\ldots} & Présenter des chiffres \tabularnewline
\en{Sales rose / fell by 12\,\%.} & Hausse / baisse \tabularnewline
\en{Compared to last year / last quarter\ldots} & Comparaison \tabularnewline
\en{Sorry to interrupt, but\ldots} & Interrompre poliment \tabularnewline
\en{Could I come in here?} & Prendre la parole \tabularnewline
\en{Can I just add something?} & Ajouter \tabularnewline
\en{Could you go back to that point?} & Revenir en arrière \tabularnewline
\en{Could you clarify what you mean by\ldots?} & Faire préciser \tabularnewline
\en{If I understand correctly, you're saying that\ldots} & Vérifier sa compréhension \tabularnewline
\en{That's a valid point, but\ldots} & Nuancer \tabularnewline
\en{I'd like to make a suggestion.} & Proposer \tabularnewline
\en{What do you think about\ldots?} & Solliciter un avis \tabularnewline
\en{Shall we take a vote?} & Voter \tabularnewline
\en{Let's put that aside for now.} & Mettre de côté \tabularnewline
\en{We're running out of time.} & Le temps presse \tabularnewline
\en{To sum up, \ldots} & Résumer \tabularnewline
\en{Any other business?} & Questions diverses \tabularnewline
\en{I'll send you the minutes.} & J'enverrai le compte rendu. \tabularnewline
\end{tabularx}
\end{tableau}

\section{Situation 16 --- S'excuser, remercier, demander un service}
\begin{tableau}[Trois moments clés de la vie sociale]
\begin{tabularx}{\linewidth}{@{}p{5.0cm}p{4.6cm}X@{}}
\entete{Situation} & \entete{Bloc} & \entete{Réponse attendue} \tabularnewline
\hline
S'excuser (léger) & \en{Sorry! / I'm sorry.} & \en{That's okay. / No problem.} \tabularnewline
S'excuser (sérieux) & \en{I'm really sorry about what happened. It was my fault.} & \en{Don't worry about it. / I understand.} \tabularnewline
S'excuser (retard) & \en{Sorry I'm late --- the traffic was terrible.} & \en{That's fine, we've just started.} \tabularnewline
S'excuser et réparer & \en{I apologise. How can I make it right?} & \en{Just\ldots\ and that'll be fine.} \tabularnewline
Remercier (simple) & \en{Thanks! / Thank you.} & \en{You're welcome.} \tabularnewline
Remercier (sincère) & \en{Thank you so much, I really appreciate it.} & \en{It was my pleasure.} \tabularnewline
Remercier (écrit) & \en{Thanks a lot for your help with the report.} & \en{Any time.} \tabularnewline
Demander un service & \en{Could you do me a favour?} & \en{Sure, what do you need?} \tabularnewline
Demander directement & \en{Would you mind helping me with this?} & \en{Of course not --- here you are.} \tabularnewline
Refuser poliment & \en{I'd love to, but I'm afraid I can't.} & \en{No problem, maybe next time.} \tabularnewline
Reporter & \en{Could we do it another time? I'm swamped today.} & \en{Sure, let me know when.} \tabularnewline
\end{tabularx}
\end{tableau}

\begin{piege}[\en{I'm sorry} a deux sens]
\en{I'm sorry} = « je m'excuse » \emph{et} « je suis désolé pour toi » (compassion).
\en{``My father died.''} --- \en{``I'm so sorry.''} est la seule bonne réponse --- jamais
\en{*It's not a problem} ou \en{*You're welcome}. Cette confusion fait commettre des
impairs énormes ; retenez : devant une mauvaise nouvelle, \en{I'm sorry} (ou
\en{That's awful, I'm so sorry for your loss}).
\end{piege}

\section{Situation 17 --- Ne pas comprendre : les phrases de survie}
\begin{memo}[Les 12 phrases qui vous sauvent dans chaque conversation]
\begin{enumerate}
  \item \en{Sorry, could you repeat that, please?} --- Pardon, pourriez-vous répéter ?
  \item \en{Could you speak more slowly, please?} --- Pourriez-vous parler plus lentement ?
  \item \en{Could you say that again more simply?} --- Pourriez-vous le redire plus simplement ?
  \item \en{What does ``\ldots'' mean?} --- Que veut dire « \ldots » ?
  \item \en{How do you say ``\ldots'' in English?} --- Comment dit-on « \ldots » en anglais ?
  \item \en{Sorry, I didn't catch that.} --- Pardon, je n'ai pas saisi.
  \item \en{Do you mean\ldots?} --- Voulez-vous dire\ldots ?
  \item \en{Let me check I understood: you're saying that\ldots} --- Je vérifie : vous dites que\ldots
  \item \en{I'm still learning English, so please bear with me.} --- J'apprends encore l'anglais, soyez indulgent.
  \item \en{Could you write it down for me?} --- Pourriez-vous me l'écrire ?
  \item \en{Could you spell that, please?} --- Pourriez-vous l'épeler ?
  \item \en{I understand most of it, but I'm not sure about\ldots} --- Je comprends l'essentiel, mais pas sûr pour\ldots
\end{enumerate}
\end{memo}

\begin{astuce}[La phrase qui détend tout le monde]
\en{My English isn't perfect yet, but I'm working on it every day.}
(Mon anglais n'est pas encore parfait, mais je travaille dessus chaque jour.)
Cette phrase, dite avec le sourire, produit trois effets : elle excuse d'avance vos erreurs,
elle montre votre sérieux, et elle pousse l'interlocuteur à ralentir et à vous aider.
Apprenez-la par c\oe ur --- elle vaut de l'or dans les six premiers mois.
\end{astuce}

\section{Les 23 autres situations à préparer vous-même}
Vous avez maintenant la mécanique. Les situations suivantes sont listées avec leurs trois
blocs essentiels. Construisez vous-même un dialogue de 8 répliques pour chacune --- c'est
l'exercice le plus rentable du livre.

\begin{tableau}[Situations 18 à 40 --- blocs de départ]
\footnotesize
\begin{longtable}{@{}p{3.4cm}p{5.6cm}p{4.8cm}@{}}
\entete{Situation} & \entete{Blocs essentiels} & \entete{Votre mission} \tabularnewline
\endfirsthead
\entete{Situation} & \entete{Blocs essentiels} & \entete{Mission} \tabularnewline
\endhead
\hline
\textbf{18. À la pharmacie} & \en{Do you have something for\ldots?} / \en{Do I need a prescription?} / \en{How often should I take it?} & Jouer la scène complète \tabularnewline
\textbf{19. À la banque} & \en{I'd like to open an account.} / \en{What's the exchange rate?} / \en{I'd like to transfer money abroad.} & Ajouter les frais et délais \tabularnewline
\textbf{20. À la poste} & \en{I'd like to send this to France.} / \en{How long will it take?} / \en{Registered, please.} & Peser, payer, confirmer \tabularnewline
\textbf{21. Chez le coiffeur} & \en{Just a trim, please.} / \en{How short would you like it?} / \en{Not too short at the back.} & Négocier la coupe \tabularnewline
\textbf{22. En taxi} & \en{Could you take me to\ldots?} / \en{How much will it be?} / \en{Could you drop me here, please?} & Discuter le prix avant \tabularnewline
\textbf{23. Louer une voiture} & \en{I'd like to rent a car for three days.} / \en{Is insurance included?} / \en{Do I need to return it with a full tank?} & Ajouter l'état des lieux \tabularnewline
\textbf{24. Acheter en ligne} & \en{When will my order arrive?} / \en{Can I change the delivery address?} / \en{I'd like a refund.} & Réclamation complète \tabularnewline
\textbf{25. Panne technique} & \en{My phone won't turn on.} / \en{It's still under warranty.} / \en{Can you fix it today?} & Décrire le problème précisément \tabularnewline
\textbf{26. À la salle de sport} & \en{I'd like to sign up.} / \en{How much is the monthly fee?} / \en{Do you have a personal trainer?} & Poser 5 questions \tabularnewline
\textbf{27. Parler football} & \en{Did you watch the match?} / \en{We won two-nil!} / \en{The referee was terrible.} & Débat de 5 minutes \tabularnewline
\textbf{28. À l'école (parent)} & \en{How is my child doing?} / \en{He needs to improve his English.} / \en{When are the exams?} & Rendez-vous complet \tabularnewline
\textbf{29. Avec un professeur} & \en{Could you explain this again?} / \en{When is the deadline?} / \en{Can I have an extension?} & Demander un délai poliment \tabularnewline
\textbf{30. Au commissariat} & \en{I'd like to report a theft.} / \en{Someone stole my phone.} / \en{I need a police report.} & Déclaration détaillée \tabularnewline
\textbf{31. Urgence} & \en{Call an ambulance!} / \en{There's been an accident.} / \en{He's not breathing.} & Message d'alerte complet \tabularnewline
\textbf{32. Inviter} & \en{Would you like to\ldots?} / \en{Are you free on Saturday?} / \en{Let's meet at 7.} & Trois invitations différentes \tabularnewline
\textbf{33. Refuser} & \en{I'd love to, but I can't.} / \en{Something came up.} / \en{Can we do it another time?} & Refus + proposition \tabularnewline
\textbf{34. Faire connaissance en ligne} & \en{How long have you been learning English?} / \en{Would you like to practise together?} / \en{Shall we do a voice call?} & Premier appel de 5 min \tabularnewline
\textbf{35. Parler de religion} & \en{I pray five times a day.} / \en{We're fasting this month.} / \en{Everyone is welcome to join us.} & Échange respectueux \tabularnewline
\textbf{36. Parler de sa famille} & \en{I come from a big family.} / \en{I have four brothers and two sisters.} / \en{My parents live in\ldots} & Description de 2 min \tabularnewline
\textbf{37. Décrire une photo} & \en{In the foreground, there's\ldots} / \en{In the background, I can see\ldots} / \en{It looks like they're\ldots} & Décrire 3 photos réelles \tabularnewline
\textbf{38. Présenter un projet} & \en{The aim of this project is to\ldots} / \en{There are three stages.} / \en{I'd be happy to take questions.} & Exposé de 3 minutes \tabularnewline
\textbf{39. Passer un appel vidéo} & \en{Can you see me? Can you hear me?} / \en{You're on mute.} / \en{Let me share my screen.} & Gérer les problèmes techniques \tabularnewline
\textbf{40. Faire un discours / toast} & \en{I'd like to say a few words.} / \en{Thank you all for coming.} / \en{Let's raise our glasses.} & Discours de 90 secondes \tabularnewline
\end{longtable}
\end{tableau}

\begin{mission}[Chapitre 14 --- les 40 jours]
Une situation par jour, pendant 40 jours. Protocole quotidien (25 minutes) :\\
\textbf{5 min} --- lire le dialogue et ses blocs, à voix haute.\\
\textbf{5 min} --- écouter (synthèse vocale ou lecture enregistrée par vous).\\
\textbf{10 min} --- jouer les deux rôles sans le livre, en changeant trois détails,
et en vous enregistrant.\\
\textbf{5 min} --- créer 5 cartes Anki avec les blocs les plus utiles du jour.\\
Au jour 40, vous posséderez environ \textbf{500 blocs de conversation actifs} --- c'est
précisément ce qui permet de tenir une conversation de 30 minutes.
\end{mission}
